\documentclass[11pt]{article}

\usepackage[margin=1in]{geometry}
\usepackage{amsmath,amssymb,amsthm,mathtools}
\usepackage{booktabs}
\usepackage{enumitem}
\usepackage{xcolor}
\usepackage[numbers,sort&compress]{natbib}
\usepackage[
  colorlinks=true,
  linkcolor=blue!45!black,
  citecolor=blue!45!black,
  urlcolor=blue!55!black,
  breaklinks=true
]{hyperref}

\newtheorem{proposition}{Proposition}
\newtheorem{corollary}{Corollary}
\newtheorem{remark}{Remark}

\newcommand{\cA}{\mathcal A}
\newcommand{\cB}{\mathcal B}
\newcommand{\cC}{\mathcal C}
\newcommand{\cD}{\mathcal D}
\newcommand{\cH}{\mathcal H}
\newcommand{\cM}{\mathcal M}
\newcommand{\cN}{\mathcal N}
\newcommand{\cP}{\mathcal P}
\newcommand{\cR}{\mathcal R}
\newcommand{\cT}{\mathcal T}
\newcommand{\id}{\operatorname{id}}
\newcommand{\Tr}{\operatorname{Tr}}
\newcommand{\Ad}{\operatorname{Ad}}
\newcommand{\dist}{\operatorname{dist}}

\title{Observer Access in Doubled DSSYK:\\
Isometric Equivalence and the Missing Resource Constraint}
\author{}
\date{July 2026}

\begin{document}

\maketitle

\begin{abstract}
We examine whether the equal-energy constraint in doubled double-scaled SYK
changes the operational information available to a relational observer.  For
identical disorder and a nondegenerate spectrum within a fixed symmetry
sector, the constrained physical Hilbert space is spectrally isometric to one
copy of DSSYK.  We show that
transporting a diary code and a sequential observer protocol through this
isometry leaves the complete observer-record channel, its distance from
diary-blind processes, and optimal recovery fidelity unchanged.  Charge-sector
and diary-twirl controls separate access to public constraint metadata from
access to private microscopic payloads.  The currently explicit kinematical
relational construction permits all bounded system operators and therefore
does not by itself supply a restricted observer protocol.  We conclude that
any nontrivial constraint-induced interpretation of the de Sitter scale as an
observer-process cutoff must derive an additional resource restriction---such
as an implementation or complexity budget---rather than follow from the
equal-energy constraint or relational dressing alone.
Recent anti-scrambling results sharpen this missing-resource statement:
their four-point observable is a proper two-fold out-of-time-order functional,
whereas a passive observer's reduced record lives on an ordinary
Schwinger--Keldysh contour.  Converting the former into the latter requires a
separately priced timefold compiler.  An explicit observer clock does not by
itself close the resource gap.  Choosing the canonical covariant time POVM for
the CLPW maximum-entropy clock gives a Cauchy one-read density matching the
tracial two-point kernel, but neither the clock state nor that POVM fixes a
post-measurement instrument.  Fresh, persistent, and contact-disturbed
instruments have the same one-read data and different two-record access.
Narovlansky--Verlinde do supply a model constraint-preserving detector
contact, but the separate observer-energy cap does not bound its normalization
or integrated action.  A recent Euclidean metric-susceptibility bound likewise
does not select a Lorentzian instrument.  Every completed clock and detector
process remains identical after isometric transport while detector action
remains unpriced.
\end{abstract}

\section{Introduction}
\label{sec:intro}

Positive cosmological constant is often associated with a finite entropy, a
bounded static patch, and a finite observer \citep{CLPW}.  These facts motivate several
inequivalent uses of the phrase ``dual cutoff'': a regulator of a continuum
dual, a bounded spectrum or finite state count, a finite radial wall, or an
operational limit on what an observer can distinguish and recover.  The last
meaning is stronger than the first three.  An entropy or an algebra type does
not, by itself, specify a channel from microscopic states to an observer's
records.

Double-scaled SYK (DSSYK) provides a sharp setting in which to separate these
claims.  Narovlansky and Verlinde proposed a doubled model with identical
disorder and an equal-energy constraint, together with constraint-preserving
dressed operators and an observer-detector interpretation
\citep{NarovlanskyVerlinde2023}.  A relational reconstruction in chord space
subsequently made the clock and gauge-invariant operator algebra explicit
\citep{AguilarRelational2025}.  Adjacent work studies thermodynamics,
correlators, operator growth, complexity, entanglement, and finite-cutoff
deformations \citep{TiettoVerlinde2025,NarovlanskyDynamics2025,
AguilarT2,AguilarFiniteCutoff}.  This raises an operational question:
does the equal-energy constraint change access to microscopic information, or
does it only change its representation and relational description?

The answer at the level of exact Hilbert-space kinematics is exact.  The
physical Hilbert space of the doubled equal-energy model is isometric to one
copy.  If states,
observer interventions, memories, records, and decoders are transported
through that isometry, the complete record channel is unchanged.  Every
record-based access quantity is consequently unchanged.  This is an
elementary invariance statement, not a new general theorem about quantum
combs.  Its significance here is diagnostic: the constraint cannot be the
source of an operational access advantage unless it also selects a protocol
class that is \emph{not} the isometric transport of the one-copy class.

The currently explicit relational construction does not make that selection.
In the chord-space kinematical description it builds conditioned observables
from arbitrary bounded system operators and gives a Type~I$_\infty$ algebra;
at finite microscopic dimension the corresponding full matrix algebra is
Type~I$_n$
\citep{AguilarRelational2025}.  It supplies relational time and dressing, but
not an implementation budget.  Thus the missing ingredient is not another
entropy or correlator.  It is a common resource theory for observer
operations on the one-copy and doubled descriptions.

Section~\ref{sec:dssyk} fixes the doubled kinematics and the distinction
between algebra and protocol.  Section~\ref{sec:invariance} states the
isometric no-free-access proposition.  Section~\ref{sec:controls} gives exact
charge and twirl controls.  Section~\ref{sec:literature} positions the result
relative to nearby DSSYK work.  Section~\ref{sec:cutoff} states what would be
required for an operational cutoff claim.  Appendix~\ref{app:shell} gives an
energy-matched finite-shell diary, and Appendix~\ref{app:comb} records the
sequential-distance convention.

\section{Equal-energy DSSYK and observer protocols}
\label{sec:dssyk}

Consider two DSSYK systems with the same disorder realization and Hamiltonians
$H_L,H_R$.  Fix a common symmetry sector in which the spectrum is
nondegenerate.  The physical-state condition is
\begin{equation}
  (H_L-H_R)|\psi\rangle_{\rm phys}=0.
  \label{eq:constraint}
\end{equation}
For a generic nondegenerate spectrum, every physical energy eigenstate is
uniquely paired,
\begin{equation}
  |E_i\rangle_{\rm phys}=|E_i\rangle_L|E_i\rangle_R,
  \qquad
  H_{\rm phys}=\frac{H_L+H_R}{2}.
  \label{eq:paired}
\end{equation}
It follows directly that
\begin{equation}
  \rho_{\rm phys}(E)=\rho_{\rm DSSYK}(E),
  \label{eq:dos}
\end{equation}
up to declared shell binning, rather than the square of the one-copy density.
The map
\begin{equation}
  W:\cH_0\longrightarrow \cH_{\rm phys},
  \qquad W|E_i\rangle=|E_i\rangle_L|E_i\rangle_R,
  \label{eq:W}
\end{equation}
is unitary onto the physical subspace and intertwines the one-copy Hamiltonian
with $H_{\rm phys}$.  Equations~\eqref{eq:constraint}--\eqref{eq:dos} are the
kinematic core of the Narovlansky--Verlinde construction
\citep{NarovlanskyVerlinde2023}.

If an energy eigenspace has multiplicity $g_E>1$, the full kernel of
$H_L-H_R$ contains $\cH_E^L\otimes\cH_E^R$ and has multiplicity $g_E^2$ at
that energy.  In that case $W$ identifies a declared diagonal paired subspace,
not the full kernel, unless an additional pairing or gauge fixing is imposed.
All statements below therefore apply sector by sector, or to any explicitly
declared paired image of $W$.

The same work constructs physical scaling operators by averaging products of
left and right operators over relative time.  One copy can be interpreted as
a clock for the other, while physical time is generated by
$H_{\rm phys}$.  Aguilar-Gutierrez derives the corresponding relational
framework in chord space.  In that kinematical description the invariant
algebra takes the schematic form
\begin{equation}
  \cA_{\rm inv}
  =\left(\cB(\cH_0^C)\otimes\cB(\cH_0^S)\right)^{\varphi},
  \qquad A\in\cB(\cH_0^S),
  \label{eq:relalg}
\end{equation}
with relationally conditioned operators obtained by group averaging
\citep{AguilarRelational2025}.  Equation~\eqref{eq:relalg} is crucial: gauge
invariance determines how an operator is dressed, but the displayed
kinematical construction allows every bounded system operator $A$.

An operator algebra and an operational protocol are not identical data.  A
protocol additionally states which operators can be implemented, at what
times and strengths, with which memories and ancillas, and under what cost
budget.  If all bounded operators and arbitrary ancillas are available, the
isometry~\eqref{eq:W} can transport the complete protocol.  If only ``simple''
operators are available, simplicity must be defined on both sides before an
access comparison is meaningful.

\section{Isometric no-free-access}
\label{sec:invariance}

Let a diary code $D$ be encoded by an isometry $V:D\to\cH_0$.  A general
$K$-step observer protocol has retained memory $M_j$ and emits records $R_j$:
\begin{equation}
  \Phi_j:
  \cH_0\otimes M_{j-1}\longrightarrow
  \cH_0\otimes M_j\otimes R_j.
  \label{eq:protocol}
\end{equation}
No Markov assumption is made: $M_j$ can retain all system-observer memory.
After the final hidden system and memory are traced out, the protocol induces
a diary-to-record channel
\begin{equation}
  \cN_K^{\cP}:D\longrightarrow R_1\cdots R_K.
  \label{eq:recordchannel}
\end{equation}

Transport the encoder and every system leg through $W$:
\begin{align}
  V'&=WV,\nonumber\\
  \Phi_j'&=
  (\Ad_W\otimes\id_{M_jR_j})\circ\Phi_j\circ
  (\Ad_{W^\dagger}\otimes\id_{M_{j-1}}),
  \label{eq:transport}
\end{align}
where $\Ad_X(\rho)=X\rho X^\dagger$.  The inverse map is needed only on the
physical image; any channel extension away from that image is operationally
irrelevant.

\begin{proposition}[Isometric no-free-access]
\label{prop:isometry}
The transported protocol has exactly the same diary-to-record channel:
\begin{equation}
  \cN_K^{\cP'}=\cN_K^{\cP}.
  \label{eq:channel-equality}
\end{equation}
If the allowed protocol class, diary-blind comparison class, and decoder class
are transported by the same isometry, then pairwise record
distinguishability, code-diamond distance from diary blindness, cumulative
sequential distance from blind comparison protocols, optimal recovery
fidelity, and unrestricted observer capacity are equal.
\end{proposition}

\begin{proof}
Insert~\eqref{eq:transport} at every step.  Each adjacent
$W^\dagger W$ cancels.  The sole remaining $W$ acts on the final hidden system
and disappears under partial trace.  This proves
\eqref{eq:channel-equality} for every diary input, including inputs entangled
with an arbitrary reference.  Trace and diamond norms are invariant under
isometric conjugation.  Transport maps blind protocols bijectively to blind
protocols and preserves every reachable-state step defect.  Taking infima
over blind comparisons therefore preserves cumulative distance.  Finally,
transporting a decoder gives the same recovered reference-diary state.
\end{proof}

For two diary states $\rho_0,\rho_1$, define
\begin{equation}
  \delta_K(\cP)
  =\frac12\left\|
    \cN_K^{\cP}(\rho_0)-\cN_K^{\cP}(\rho_1)
  \right\|_1.
  \label{eq:delta}
\end{equation}
A diary-blind record channel is constant on the declared code.  For every such
channel $\cC_K$, the triangle inequality gives, using the full diamond norm,
\begin{equation}
  \left\|\cN_K^{\cP}-\cC_K\right\|_{\diamond,D}
  \geq \delta_K(\cP).
  \label{eq:blindlower}
\end{equation}
Both sides of~\eqref{eq:blindlower} are invariant under the transport above.

\begin{corollary}[Equal-energy DSSYK]
\label{cor:dssyk}
The equal-energy constraint, paired representation, and one-copy physical
density of states cannot by themselves change diary distinguishability or
recovery capacity.  The exact one-copy observer operation
\begin{equation}
  A_{\rm one}=W^\dagger A_{\rm phys}W
  \label{eq:operatortransport}
\end{equation}
reproduces the complete record channel of the physical doubled operation.
Any operational difference requires a protocol restriction not related by
this transport.
\end{corollary}

The corollary does not say that every transported operator is easy to
implement.  It says precisely that implementation cost is the additional
physical datum needed to make the question nontrivial.

\section{Exact controls: metadata and blindness}
\label{sec:controls}

Two finite controls calibrate what the proposition does and does not exclude.

\subsection{Charge metadata is not fixed-sector payload}

Let
\begin{equation}
  \cH=\bigoplus_q\cH_q
  \label{eq:charge-decomp}
\end{equation}
and consider an observer that records only the sector label,
\begin{equation}
  \cN_Q(\rho)
  =\sum_q\Tr(P_q\rho)|q\rangle\langle q|_R.
  \label{eq:charge-channel}
\end{equation}
States in different sectors are perfectly distinguishable:
\begin{equation}
  \delta_Q(\rho_q,\rho_{q'})=1\qquad(q\neq q').
  \label{eq:header}
\end{equation}
Within one sector, however,
\begin{equation}
  \cN_Q(\rho)=|q\rangle\langle q|_R
  \quad\text{for all }\rho\text{ supported in }\cH_q.
  \label{eq:payloadblind}
\end{equation}
Thus fixed-sector pairwise distinguishability is zero.  If
$d_q=\dim\cH_q$ and half of a maximally entangled payload state is sent
through~\eqref{eq:charge-channel}, every decoder of the constant record has
entanglement fidelity at most
\begin{equation}
  F_e\leq \frac{1}{d_q^2}.
  \label{eq:Fe}
\end{equation}
We use the convention in which fidelity with a pure target is its overlap.
The control exposes a constraint label perfectly while carrying no private
payload inside a fixed sector.  Symmetry-constrained recovery can of course
show richer behavior when the physical channel acts nontrivially within
sectors \citep{NakataClouded}; the point of
\eqref{eq:charge-channel} is to calibrate the metadata-only endpoint.

\subsection{Same-shell diary twirl}

For a $d$-dimensional code, let
\begin{equation}
  \cT_D(\rho)=\int dU\,U\rho U^\dagger
  =\frac{I_D}{d}\Tr\rho,
  \qquad
  \cC_K^{\rm twirl}=\cN_K\circ\cT_D.
  \label{eq:twirl}
\end{equation}
The control uses the same shell and observer instrument as the actual channel
but is exactly diary blind, so $\delta_K(\cC_K^{\rm twirl})=0$.

For a binary classical diary with outputs $\sigma_0,\sigma_1$, the uniform
twirl outputs $\bar\sigma=(\sigma_0+\sigma_1)/2$.  Directly,
\begin{equation}
  \left\|\cN_K^{\rm cq}-\cC_K^{\rm twirl}\right\|_\diamond
  =\frac12\|\sigma_0-\sigma_1\|_1
  =\delta_K.
  \label{eq:twirlidentity}
\end{equation}
The twirl therefore saturates the pairwise lower bound
\eqref{eq:blindlower} for the binary cq restriction.  For a coherent quantum
code, reference-assisted inputs may give a larger channel distance; the
universal statement remains~\eqref{eq:blindlower}.

The controls establish a three-way distinction:
\begin{enumerate}[label=(\roman*)]
  \item \emph{metadata access}: public energy, charge, or clock labels;
  \item \emph{directional access}: selected microscopic diary directions;
  \item \emph{payload access}: controlled recovery of a declared code.
\end{enumerate}
Nonzero correlators, nontrivial dressing, or a changed algebra type do not by
themselves determine which branch is realized.

\section{Relation to existing DSSYK results}
\label{sec:literature}

The doubled equal-energy construction and its simple dressed detector belong
to Narovlansky and Verlinde \citep{NarovlanskyVerlinde2023}.  The
Rahman--Susskind line is an independent DSSYK/de Sitter proposal and a
competing separation-of-scales interpretation, not a co-origin of the
equal-energy model
\citep{SusskindSeparation,RahmanDSJT,RahmanSusskindComments}.  This distinction matters
because our conclusion concerns the doubled constraint, not the validity of
the broader DSSYK/de Sitter program.

The adjacent literature supplies several parts of the story without the
record-channel comparison made here.  Relational chord-space work constructs
the constraint sector, clock states, and invariant algebra
\citep{AguilarRelational2025}.  Observer and finite-cutoff papers compute
entropy, thermodynamics, spectra, correlators, OTOCs, Krylov complexity, and
algebraic entanglement
\citep{TiettoVerlinde2025,NarovlanskyDynamics2025,AguilarT2,
AguilarFiniteCutoff}.  Algebraic work characterizes the double-scaled
operator algebras \citep{XuAlgebras}; recent $q$-Askey deformations further
show that the resulting factor type depends on which operators are included
\citep{AguilarQAskey}.  Rajgadia and Xu provide the closest
access-language near hit: state-adapted dressed operators restore sensitivity
to the purity of a selected Kourkoulou--Maldacena state and change the
emergent algebra type \citep{RajgadiaXu}.  That construction does not define
an equal-energy shell diary, a memoryful record channel, or an isospectral
recovery comparison.  Single-sided DSSYK reconstruction and commutant results
are likewise algebraic rather than diary-channel calculations
\citep{CaoGao}.

Several contemporaneous papers on observer correlators give the closest
dynamical comparison.  Observer-dressed de Sitter OTOCs exhibit time advances
and, at tree level, maximal Lyapunov growth
\citep{MilekhinNarovlanskyXu2026}.  Including observer recoil
and gravitational backreaction obstructs a naive tracial interpretation
through failures of cyclicity or positivity
\citep{ChenStanfordTangYang2026,CuiKolchmeyer2026}.  Harlow and Zhao show that
the milder anti-scrambling correlator can instead be represented in a bounded
quantum system by using folded Euclidean evolution, while emphasizing that no
general prescription for all correlators is yet selected
\citep{HarlowZhao2026}.  These results make the distinction at issue here more
acute: a worldline correlator, even one with a physically motivated dressing
or time contour, is not yet a sequential diary-to-record channel.  In
particular, the ordinary forward-time observer evolution discussed in
Ref.~\citep{CuiKolchmeyer2026} is insensitive to the relevant OTOCs.  The
cluster does not supply an equal-energy shell diary, a recovery map, or the
isospectral one-copy control used below.  Conversely, its backreaction
results caution against promoting a kinematical tracial or Type~II structure
to a nonperturbative observer-process claim.

Positive higher-order operational representations of OTO experiments are
also known.  The out-of-time-order tensor packages forward and reversed
evolution into a quantum process with an information-theoretic interpretation
\citep{ZonniosEtAl2021}.  It therefore precludes claiming novelty for the
broad move from an OTOC to a process.  It does not derive the availability or
cost of reversal, copies, or coherent control from a de Sitter observer
construction, which is the narrower issue here.

A bounded primary-source scan of these works found no computation combining a
declared microscopic diary code, a sequential observer record, a distance
from diary blindness or recovery fidelity, and an isospectral equal-energy
control.  We do not claim that constraints, relational observables, quantum
combs, or symmetry-modified recovery are individually new.  Nor do we claim
an exhaustive no-prior-art result beyond the checked DSSYK neighborhood.
Information recovery has been studied in de Sitter settings by other methods,
including islands and observer backreaction \citep{AalsmaSybesma}; those
works answer a different channel question.

\section{The missing resource constraint}
\label{sec:cutoff}

Corollary~\ref{cor:dssyk} closes the unrestricted mechanism
\begin{equation}
  \begin{aligned}
    &\text{equal-energy constraint}
      +\text{paired representation}
      +\text{full relational algebra}
      \\
    &\hspace{3cm}\Longrightarrow\quad
      \text{new operational access}.
  \end{aligned}
  \label{eq:deadmechanism}
\end{equation}
The implication in~\eqref{eq:deadmechanism} is false: exact transport gives an
identical record channel.

This null is comparative and specific to the equal-energy constraint.  It
does not say that the compact one-copy DSSYK spectrum, finite time resolution,
or $q$-deformation cannot define an intrinsic bandwidth or cutoff.  Any such
limitation is shared by the isospectral one-copy and paired descriptions and
must be analyzed as a property of DSSYK itself, rather than as access created
by the equal-energy constraint.

A narrower question remains.  The constraint may change the \emph{cost} of
implementing a code-sensitive record.  To state that question without
choosing the answer, let $V:D\to\cH_0$ be an encoder, let
$\cR:R_1\cdots R_K\to D$ be a decoder, and introduce a common implementation
cost $c(V,\cP,\cR)$.  Define
\begin{equation}
  C_{\rm obs}(K,\epsilon;B)
  =\sup_{D,V,\cP,\cR}
    \left\{\log\dim D:\begin{array}{l}
      c(V,\cP,\cR)\leq B,\\[-1mm]
      1-F_e\!\left(\cR\circ\cN_{K,V}^{\cP}\right)\leq\epsilon
    \end{array}\right\}.
  \label{eq:Cobs}
\end{equation}
Here $\cN_{K,V}^{\cP}$ includes the encoder, and $F_e$ is entanglement fidelity
on the declared diary code.  A convention that treats encoding or decoding as
free is allowed only if stated on both sides.  The cost must be defined on the
one-copy and doubled descriptions before the diary calculation, and the two
allowed classes must be matched by the same physical rule.  Candidate
resources include word length in a declared set of simple matter/chord
generators, allowed numbers and scaling dimensions of insertions, integrated
detector coupling and contact count, or operator/Krylov complexity relative to
a fixed generator set.  None of these detector-level costs is selected here.
The observer construction does motivate a lower-bounded clock Hamiltonian and
particular clock states.  We evaluate the canonical covariant-time completion
below, while keeping separate the instrument that the bulk construction does
not select.  Complexity and finite-cutoff
observables already have a rich DSSYK literature, so adopting one of them as
an access cost would require an additional operational derivation rather than
a relabeling of an existing curve.

Recent anti-scrambling work supplies two physically motivated structures that
could enter such a derivation: the observer-energy endpoint produces a
time-smeared two-point function, and a bounded Hamiltonian permits backwards
segments of a folded Euclidean correlator
\citep{ChenStanfordTangYang2026,HarlowZhao2026}.  Neither fact alone fixes a
deterministic observer comb.  The first is low-point transfer data with
multiple possible process completions; the second is a nonunitary operation
if interpreted as direct state evolution.  Turning either into access
therefore still requires an explicit record protocol and a resource audit.

\subsection*{A clock state fixes a one-read kernel, not an instrument}
The observer algebra of Ref.~\citep{CLPW} contains a clock with
$H_{\rm obs}=q\geq0$.  For a normalized wavefunction
$f\in L^2(\mathbb R_+,dq)$, the canonical covariant time POVM has generalized
kets and effects
\begin{equation}
 |s)=\frac{1}{\sqrt{2\pi}}\int_0^\infty dq\,e^{-iqs}|q\rangle,
 \qquad E(ds)=|s)(s|\,ds .
 \label{eq:clockpovm}
\end{equation}
Such half-line covariant observables admit Naimark dilations, but a POVM does
not specify the retained dilation or post-measurement state
\citep{EgusquizaMuga1999,LeppajarviSedlak2020}.  Its one-read density is
$p_f(s)=|(s|f)|^2$.  The characteristic function is the exact half-line
autocorrelation
\begin{equation}
 M_f(\omega)
 =\int_{\max(0,-\omega)}^\infty dq\,f(q+\omega)f(q)^*
 =\int ds\,p_f(s)e^{i\omega s}.
 \label{eq:clockoverlap}
\end{equation}

The maximum-entropy wavefunction used by CLPW is
$f_\beta(q)=\sqrt\beta e^{-\beta q/2}$.  For the canonical POVM,
\begin{equation}
 p_\beta(s)=\frac{\beta/2}{\pi[(\beta/2)^2+s^2]},
 \qquad M_{f_\beta}(\omega)=e^{-\beta|\omega|/2}.
 \label{eq:clockcauchy}
\end{equation}
With $\beta_{\rm dS}=2\pi R_{\rm dS}$ this matches the transfer factor in the
tracial two-point function of Ref.~\citep{ChenStanfordTangYang2026}.  The
origin of that factor there is the observer-energy endpoint integration.  The
equality therefore identifies a canonical positive completion; it does not
derive a measurement channel from the correlator.  Moreover, CLPW explicitly
note that the maximum-entropy clock is not semiclassical: its clock-time
uncertainty is $O(\beta_{\rm dS})$.  Their semiclassical family instead has
$f_\epsilon(q)=\sqrt\epsilon g(\epsilon q)$ and
\begin{equation}
 M_{f_\epsilon}(\omega)
 =\int_{\max(0,-\epsilon\omega)}^\infty du\,
 g(u+\epsilon\omega)g(u)^*,
 \label{eq:semiclassicaloverlap}
\end{equation}
with bandwidth $O(1/\epsilon)$.

To turn $p_f$ into a system channel one may choose the classical random-offset
completion
\begin{equation}
 \Phi_f^H(\rho)=\int ds\,p_f(s)e^{-iHs}\rho e^{iHs}.
 \label{eq:clockchannel}
\end{equation}
This is completely positive, but it is additional instrument data.  Indeed,
for every normalized family $\{\sigma_s\}$,
\begin{equation}
 \mathfrak I^\sigma(ds)(\rho_C)
 =\Tr[E(ds)\rho_C]\,\sigma_s
 \label{eq:clockinstrument}
\end{equation}
is a valid measure-and-prepare instrument with the same POVM.  After a clock
interval $\tau$, its two-read law contains
$\Tr[E(ds_2)e^{-iq\tau}\sigma_{s_1}e^{iq\tau}]$ and therefore depends on
$\sigma_{s_1}$.

\begin{proposition}[Clock-state/instrument nonidentifiability]
\label{prop:clockinstrument}
A lower-bounded clock Hamiltonian, an initial clock state, and a covariant
time POVM determine the one-read density and overlap~\eqref{eq:clockoverlap}.
They do not determine a two-slot clock process.  Compatible instruments can
have identical one-read statistics and different observer-record channels.
\end{proposition}

Freshly measuring, discarding, and preparing $|f\rangle$ gives independent
offsets.  Retaining the first Naimark pointer and reusing it without contact
disturbance gives a common offset.  A detector-induced pointer kick gives a
third process.  None is selected by Ref.~\citep{CLPW}, which does not specify a
POVM, instrument, or observer--field contact.

The distinction remains operational after adding a finite detector.  Take two
energy bins and the temporal diary
\begin{equation}
 H_j=\Delta E|1\rangle\!\langle1|_j,
 \qquad
 |\chi_\pm\rangle=\frac{|01\rangle\pm|10\rangle}{\sqrt2}.
 \label{eq:relationaldiary}
\end{equation}
A common offset preserves the exchange coherence, while fresh offsets
multiply it by $|M_f(\Delta E)|^2$.  Let one memory qubit $R$, initially in
$|0\rangle$, meet each bin through the finite contacts
\begin{equation}
 V_j=P_{+,j}\otimes I_R+P_{-,j}\otimes X_R,
 \qquad P_{\pm,j}=\frac{I\pm X_j}{2}.
 \label{eq:finitecontacts}
\end{equation}
The ordered unitary $V_2V_1$ records $X_1X_2$ parity.  Its two-record diary
distance is exactly
\begin{equation}
 \delta_{2,{\rm common}}=1,
 \qquad
 \delta_{2,{\rm fresh}}=|M_f(\Delta E)|^2.
 \label{eq:instrumentdistances}
\end{equation}
An unobserved symmetric pointer kick $\pm\kappa$ between contacts instead
gives the fixed parity record distance $|\cos(\kappa\Delta E)|$.  These
positive causal combs have the same pre-contact one-read density and include
the contact backaction in an explicit unitary dilation; the one-read clock
data choose none of them.

Every completion is also an exact comparative null.  Extending the equal-energy
isometry by the identities on clock and memory and transporting
Eq.~\eqref{eq:finitecontacts} gives
\begin{equation}
 \delta_K^{\rm doubled}=\delta_K^{\rm one\ copy}
 \label{eq:clocknull}
\end{equation}
for each fresh, common, or contact-disturbed instrument, with the same equality
for blind-comb distance and recovery fidelity.  Supplying an instrument makes
the operational assumption explicit; it does not create a constraint
advantage.

The clock calculation therefore has a narrower cutoff interpretation.  The
maximum-entropy state supplies a radius-scale soft kernel only after the
canonical completion is chosen.  Semiclassical clocks purchase bandwidth
$O(1/\epsilon)$.  If the model-specific observer-energy cap of
Ref.~\citep{TiettoVerlinde2025} is additionally imposed, its formal endpoint
again gives a Planckian hard scale with the de Sitter radius canceled.  None
of these data bounds $\int dt\,\|H_{\rm int}(t)\|$, contact count, allowed
scaling dimensions, or implementation error.  Current bulk constructions
therefore do not yet determine an operational observer cutoff.

\subsection*{Why the available detector and backreaction inputs do not compose}
It is too strong to say that the doubled proposal contains no observer--field
contact.  Narovlansky and Verlinde couple a model Unruh detector to their
constraint-preserving dressed operators through
\begin{equation}
 S_{\rm int}=\int d\tau\,
 \left[X^+(\tau)O_\Delta^-(\tau)
       +X^-(\tau)O_\Delta^+(\tau)\right],
 \qquad [H_L-H_R,O_\Delta^\pm]=0,
 \label{eq:nvdetector}
\end{equation}
and obtain
$\dot P(E_i\!\to E_j)=|X_{ij}|^2G_\Delta(E_j-E_i)$
\citep{NarovlanskyVerlinde2023}.  This gives a physically motivated form for
a proper-time contact and a response function.  It does not fix the matrix
elements $X_{ij}$, switching profile, probe readout, retained memory, or
number of contacts.

Standard localized-probe measurement theory can turn a specified coupling,
probe state, and readout into a completely positive instrument, including
causal composition of several coupling regions
\citep{FewsterVerch2020,PoloGomezGarayMartinMartinez2022}.  Influence-functional
methods can likewise represent specified open-system dynamics as a multitime
process tensor \citep{JorgensenPollock2019}.  These constructions provide the
right formal machinery for completing Eq.~\eqref{eq:nvdetector}; they take the
coupling, preparation, and readout as input and therefore do not select their
gravitational cost.

The observer thermodynamics does not supply that cost indirectly.  In a fixed
deficit-angle sector, Ref.~\citep{TiettoVerlinde2025} imposes
\begin{equation}
 0\leq E_{\rm obs}\leq E(\psi),
 \label{eq:observerenergycap}
\end{equation}
but does not relate this free-observer energy cap to the normalization or
duration of Eq.~\eqref{eq:nvdetector}.  The logical gap can be isolated as a
simple proposition.

\begin{proposition}[Energy-cap/action-budget nonimplication]
\label{prop:energycap}
Fix a free detector Hamiltonian and a cap
$E_{\rm obs}\leq E_{\max}$.  Suppose the allowed models contain a nonzero
constraint-preserving diary-sensitive interaction difference
\begin{equation}
 \Delta H(t)=H(t)-H^{(0)}(t)
 =g(t)X\otimes\Delta O_{\rm phys}(t),
 \qquad [H_L-H_R,\Delta O_{\rm phys}(t)]=0,
 \label{eq:resourceinteraction}
\end{equation}
and remain admissible under positive rescalings of its overall normalization
because the same physical construction imposes no normalization bound.
Then the energy cap alone implies no finite upper bound on
\begin{equation}
 G_D(T)=\int_0^Tdt\,\|\Delta H(t)\|.
 \label{eq:GDcap}
\end{equation}
\end{proposition}

\begin{proof}
For any $c>0$, replacing $g(t)$ by $cg(t)$ leaves the free detector
Hamiltonian, its spectrum, the cap, and constraint preservation unchanged.
It sends $G_D(T)$ to $cG_D(T)$ and sends the detector transition rates in
Eq.~\eqref{eq:nvdetector} to $c^2$ times their former values.  Hence no finite
function of $E_{\max}$ alone bounds $G_D(T)$.  Even if additional physics
later bounds the interaction pointwise, unrestricted duration or contact
count can still make the integrated action grow.
\end{proof}

The closest bounded backreaction near-hit is a recent Euclidean Gaussian
clock-detector with a smeared metric source.  It obeys the sufficient
stable-channel condition
\begin{equation}
 \frac{\Lambda_{\rm clk}^2}{\Omega_0^2}<\delta_P
 \label{eq:euclideanadmissibility}
\end{equation}
for controlled quadratic backreaction on a chosen projected gravitational
channel \citep{EspindolaAli2026}.  This is a genuine metric-susceptibility
bound, but its source strength, smearing, and detector gap are model inputs;
the construction is Euclidean and does not provide a Lorentzian readout,
matter-diary contact, or relation between
Eq.~\eqref{eq:euclideanadmissibility} and $G_D(T)$.  Thus the explicit DSSYK
contact, observer-energy cap, and Euclidean admissibility inequality are
separately useful but do not compose into the resource law required by
Eq.~\eqref{eq:Cobs}.

Transporting Eq.~\eqref{eq:nvdetector} through the equal-energy isometry also
preserves its norm, detector response, every instrument built from the same
probe preparation and readout, and $G_D(T)$.  A nontrivial doubled advantage
would still require an independently derived locality or implementation rule
that is not merely the transported one-copy rule.

\subsection*{A concrete missing resource: contour depth}
The standard chaos four-point function is a proper 2-OTO functional: its
alternating early/late ordering requires two forward and two backward contour
legs \citep{HaehlLoganayagamNarayanRangamani2017}.  This gives a precise
version of the mismatch between the recent gravitational OTOCs and an
observer record.

\begin{proposition}[Passive-record contour bound]
\label{prop:contour}
Let a detector initially uncorrelated with the static-patch fields interact
through forward-time unitary couplings, with causal adaptive operations
dilated into a detector memory and one final POVM.  Every linear statistic of
that final record depends only on 1-OTO field functionals.  A proper 2-OTO
four-point functional is not such a statistic unless the protocol class is
enlarged.
\end{proposition}

\begin{proof}
Write the joint propagator as
$U=\mathcal T\exp[-i\int dt\,H_{\rm int}(t)]$.  A final record probability is
linear in
$\Tr_S[U(\rho_D\otimes\rho_S)U^\dagger]$.  The Dyson expansions of $U$ and
$U^\dagger$ contain field correlators of the form
\begin{equation}
 \Tr_S\!\left[
 \widetilde{\mathcal T}\{\phi(t'_1)\cdots\phi(t'_m)\}
 \rho_S
 \mathcal T\{\phi(t_1)\cdots\phi(t_n)\}
 \right],
\end{equation}
which lie on one closed-time Schwinger--Keldysh contour.  Deferred
measurement absorbs causal adaptivity into a larger forward-evolving memory
without adding a contour fold.  Proper 2-OTO orderings are absent.
\end{proof}

This is also the model-specific conclusion of Cui and Kolchmeyer: for an
ordinary observer--field coupling, the reduced observer density matrix is
insensitive to the anti-scrambling OTOCs
\citep{CuiKolchmeyer2026}.  It does not mean that OTOCs are unmeasurable.
They can be reconstructed after adding Hamiltonian reversal, coherent
ancillas or copies, randomized measurements, repeated global state
preparations, or a measurement of the probe's own OTO correlator
\citep{ChaudhuriLoganayagam2018,BlocherEtAl2020,ZonniosEtAl2021}.  Those are different
resource declarations.  In particular, the forward-only reconstruction for
an arbitrary operator in Ref.~\citep{BlocherEtAl2020} uses a basis of prepared
states and superpositions and up to $2d^2$ expectation values; it is not one
passive sequential record.

The Euclidean-fold proposal is another such compiler.  It introduces a
backwards imaginary-time segment, depends on the chosen operator ordering and
an undetermined integer fold, and is not yet a general correlator rule
\citep{HarlowZhao2026}.  Directly implementing that nonunitary segment is a
heralded operation rather than a deterministic channel.  Thus the
anti-scrambling result identifies the missing resource more sharply---contour
conversion---but does not itself supply the positive diary-to-record tester
needed in~\eqref{eq:Cobs}.

\subsection*{Objection: scaling operators are physically selected}
The strongest apparent counterexample is that the Narovlansky--Verlinde
scaling operators are already physically selected: they are proposed as the
operators dual to massive bulk fields \citep{NarovlanskyVerlinde2023}.  We
grant that this supplies physical motivation for an operator family.  It does
not yet supply the common resource rule required by~\eqref{eq:Cobs}.
Selection answers which dressed observables the bulk dictionary singles out;
a resource rule must additionally say which dimensions, compositions,
couplings, time supports, and accuracies are affordable, using the same
budgeting principle in both descriptions.  Moreover, every selected doubled
operator has the transported one-copy representative $W^\dagger O W$, with
identical record-channel performance under the transported protocol class.
Consequently, diary sensitivity of a scaling operator would demonstrate
physically motivated directional sensitivity.  It would demonstrate an
access advantage caused by the constraint only if the bulk interpretation
also derived a resource budget that distinguishes the two implementations
and is not merely imposed after observing the signal.

This identifies the precise burden behind an operational dual-cutoff claim.
A bounded spectrum, finite entropy, or Type~II limit may motivate a finite
observer \citep{CLPW,DeVuystObserver}, but it does not fix
$c(V,\cP,\cR)$.  Conversely, if a microscopic limit
derives a budget $B(\Lambda)$ and one finds nontrivial scaling of
$C_{\rm obs}(K,\epsilon;B(\Lambda))$, then the de Sitter scale would control an
observer-relative access bandwidth.  That would be a defensible operational
meaning of ``dual cutoff.''  It would still not imply that $\Lambda$ is a
literal UV regulator of a fundamental dual theory.

The native spectrum gives a useful check on this distinction.  In the
Narovlansky--Verlinde normalization,
$E(\theta)=-2\mathcal J\cos\theta/\lambda$ for
$0\leq\theta\leq\pi$, so its width is
$B_{\rm nat}=4\mathcal J/\lambda$.  The state-independent bandwidth speed
limit is therefore $t_\perp\geq\pi/B_{\rm nat}$
\citep{NessBandwidth}.  Using $\mathcal J=1/R_{\rm dS}$ and
$R_{\rm dS}/G_N=8\pi/\lambda$ gives
\begin{equation}
  B_{\rm nat}=\frac{1}{2\pi G_N},
  \qquad t_\perp\geq 2\pi^2G_N.
  \label{eq:nativebandwidth}
\end{equation}
Thus the full-band native time is microscopic in this dictionary: the de
Sitter radius cancels.  A $\Lambda$-dependent observer bandwidth would have to
enter through a physically selected shell or detector/control budget, not
from boundedness of the full spectrum alone.  Equation~\eqref{eq:nativebandwidth}
is an orthogonalization bound, not a record or recovery bound.

The kinematical conclusion also leaves open the possibility that a
restricted algebra or cost emerges in a controlled thermodynamic,
semiclassical, or observer-backreaction limit.  What it excludes is deriving
that restriction from equal-energy pairing alone.  Until such a resource
class is obtained, computing a nonzero diary signal for a chosen dressed
operator would establish sensitivity of that operator family, not an access
advantage caused by the cosmological constraint.

\section{Discussion}
\label{sec:discussion}

The result is a demarcation rather than a dynamical solution.  The
equal-energy constraint makes the doubled description relational and selects
gauge-invariant dressings, but its physical state space remains isometric to
one DSSYK.  Exact protocol transport therefore preserves every observer
record and recovery task.  The charge and twirl controls show why public
constraint labels, selected microscopic sensitivity, and recoverable payload
must be tested separately.  The clock calculation adds an exact general
one-read overlap and explicit inequivalent process completions.  It shows that
the missing datum begins one step earlier than detector cost: the bulk input
does not yet select a clock instrument.  Every chosen instrument and detector
contact nevertheless transports to the one-copy control unchanged.

The conclusion can be stated compactly:
\begin{quote}
The doubled equal-energy constraint cannot create information access under
isometrically transported protocol classes.  Every access difference in the
explicit construction is therefore a claim about a restriction on observer
resources.  A chosen canonical clock POVM fixes a one-read kernel, but the
current bulk dictionary does not yet derive the clock instrument or detector
restriction needed for a constraint-induced difference.
\end{quote}
The first sentence is exact under the transported-class hypotheses.  The
second is bounded to the current bulk constructions reviewed above.
The detector audit makes that scope more precise: a model worldline contact
exists, but the checked observer-energy and Euclidean backreaction bounds do
not price its Lorentzian use.
This is compatible with DSSYK remaining useful for de Sitter holography.  It
locates the additional input that an operational interpretation requires.

\appendix

\section{An energy-matched finite-shell diary}
\label{app:shell}

The no-free-access result does not require a particular diary.  Nevertheless,
an explicit finite-shell code shows that the access question can be posed
without changing the energy header.

For each finite-$N$ disorder realization $\omega$, choose the $2m$ paired
eigenvalues nearest a target energy $E_*$ and let
\begin{equation}
  P_\omega
  =\sum_{a=1}^{2m}|E_a,E_a\rangle\langle E_a,E_a|.
  \label{eq:shellprojector}
\end{equation}
The induced width $\Delta E_\omega$ is reported per realization.  A useful
microcanonical regime requires the mean level spacing to be much smaller than
$\Delta E_\omega$, which in turn must be smaller than the scale on which the
spectral and operator envelopes vary.

Define a binary classical phase diary
\begin{align}
  |d_0\rangle_\omega
    &=\frac{1}{\sqrt{2m}}\sum_{a=1}^{2m}|E_a,E_a\rangle,
    \nonumber\\
  |d_1\rangle_\omega
    &=\frac{1}{\sqrt{2m}}\sum_{a=1}^{2m}(-1)^a
      |E_a,E_a\rangle.
  \label{eq:phasecode}
\end{align}
The states are orthogonal and have identical fine-grained energy
probabilities.  Their distinction is a phase pattern.  Ordering by energy
rank defines the code before observer outcomes and preserves its label per
realization.  The first calculation, if a resource class is later derived,
should remain per realization and report the mean and variance of
$\delta_K(\omega)$; averaging the channel before inserting the code can erase
the question by definition.

Equation~\eqref{eq:phasecode} is a binary classical code, not yet a protected
logical qubit.  A quantum payload extension must additionally control
coherences against the declared public energy algebra.

\section{Sequential blind-process distance}
\label{app:comb}

For completeness, compare the actual step maps $\Phi_j$ with maps $\Psi_j$
whose composed record channel is constant on the diary code.  Let $\eta_j$ be
the supremum trace-norm difference between the two step outputs on all states
reachable by the hybrid processes that use either actual or comparison maps
in earlier slots.  The standard telescoping argument for sequential quantum
strategies \citep{GutoskiWatrous,GutoskiStrategy} gives
\begin{equation}
  \|\cN_K-\cC_K\|_{\diamond,D}
  \leq \sum_{j=1}^K\eta_j.
  \label{eq:hybrid}
\end{equation}
For time-binned unitary interactions generated by $H_j(t)$ and diary-blind
comparisons generated by $H_j^{(0)}(t)$ on a common reachable sector,
time-ordered Duhamel evolution gives
\begin{equation}
  \|\cN_K-\cC_K\|_{\diamond,D}
  \leq 2G_D,
  \qquad
  G_D=\sum_j\int_{I_j}\!dt\,
  \|H_j(t)-H_j^{(0)}(t)\|.
  \label{eq:detectoraction}
\end{equation}
Hence two orthogonal diary states recoverable with trace-norm error at most
$\alpha$ require $G_D\geq(1-\alpha)/2$.  Large diary-blind native evolution
does not enter this necessary condition; the resource is integrated
diary-sensitive detector action.
Define
\begin{equation}
  A_K=\inf_{\substack{(\Psi_j):\\
                      \cC_K\,\mathrm{blind}}}
      \sum_{j=1}^K\eta_j.
  \label{eq:AK}
\end{equation}
Reliable recovery requires an order-one $A_K$, but order-one $A_K$ is not
sufficient: a channel may expose one charge or one diary direction without
supporting generic recovery.  Proposition~\ref{prop:isometry} preserves
$A_K$ because the isometry bijects the reachable sets and blind comparison
protocols and leaves every $\eta_j$ unchanged.

The full trace and diamond norms used here range from zero to two.  With a
normalized half-diamond convention, the right side of
\eqref{eq:blindlower} would be $\delta_K/2$.

\begingroup
\raggedright
\bibliographystyle{unsrtnat}
\bibliography{refs}
\endgroup

\end{document}
